\chapter{Discussion and Limitations}
\label{ch:discussion}

Chapter~\ref{ch:results} showed that the evidence is uneven across the five
research questions. The internal model comparison, historical device
measurements and stationary two-node workflow produced quantitative
observations. A 30-minute complete-stack run established bounded concurrent
service survival while exposing unmet Scout latency and strict swap criteria. Comparative
temporal-filter effects, formal mission reliability and external transfer were
not established.
This chapter interprets that pattern rather than treating every implemented
component as an experimentally validated contribution.

The central hypothesis proposed that a resource-aware tiered edge
architecture can divide hazard-detection responsibilities between heterogeneous
drones. The results support the narrower claim that the implemented Sentinel
and Scout can execute that division under bounded stationary conditions and
that the division exposes measurable predictive, runtime and latency
trade-offs. They do not establish that the tiered workflow improves accuracy in general,
that the selected pair is optimal, or that the system is ready for autonomous
industrial use.

\section{Answers to the research questions}
\label{sec:discussion:1}

Table~\ref{tab:discussion-rq-answers} gives the final evidence-bounded answer to
each research question. ``Bounded answer'' means that the retained evidence
answers part of the question under stated conditions while another evidence
layer remains unavailable. ``Not evaluated'' means that no eligible experiment
supports a performance answer.

\begin{table}[H]
  \centering
  \caption{Evidence-bounded answers to the five research questions.}
  \label{tab:discussion-rq-answers}
  \scriptsize
  \begin{tabularx}{\textwidth}{@{}p{0.07\textwidth}p{0.17\textwidth}X@{}}
    \toprule
    RQ & Status & Answer supported by the accepted evidence \\
    \midrule
    RQ1 & Bounded answer & The six candidates exhibited different
    predictive and laptop-runtime trade-offs, and the historical Sentinel and
    Scout pair executed on their devices. The strengthened 12-cell screen was
    not completed, so the experiment does not select a supported candidate for
    every platform. \\
    RQ2 & Bounded answer & A stationary 30-minute complete-stack run
    completed without process death, OOM, thermal shutdown or logger loss while
    both flight-service paths were active. It missed the Scout p95 target and
    showed 0.922~MiB Sentinel swap growth, so not every acceptance criterion was
    satisfied. \\
    RQ3 & Bounded answer & The Scout produced both correlated confirmations
    and rejections, and the added decision time was measured. The effects of
    the temporal filter and tiered workflow on event retention and false-alert burden
    are not estimable from the retained frame-level evidence. \\
    RQ4 & Bounded answer & Two nominal branches, four selected fault paths
    and one hash-linked visual confirmed-hazard mission completed in development
    simulation. The formal repeated campaign and physical two-vehicle mission
    were not completed, so correct behaviour is demonstrated only for the
    retained pilots. \\
    RQ5 & Not evaluated & Neither external source reached a frozen eligible
    inference manifest. No transfer-performance conclusion is supported. \\
    \bottomrule
  \end{tabularx}
\end{table}

\subsection*{RQ1: predictive performance and platform support}

The internal comparison rejects the idea of a universal detector winner.
YOLOv8n-416 led validation AP50--95 and F1, while EfficientDet led test
AP50--95, precision and F1. All candidates lost F1 between validation and test,
and their grouped intervals overlapped. The ranking reversal is compatible
with sensitivity to split composition, source domains, target scales and model
inductive biases. It does not justify replacing the validation-frozen selection
with the model that happened to rank first on the test split.

Runtime also changed the ordering. EfficientDet's stronger test result used an
available CPU/XNNPACK path whose laptop inference call was approximately two
orders of magnitude slower than the GPU-capable YOLO paths. Conversely,
YOLOv8n-640 had the lowest observed YOLO laptop time but lower validation AP
than YOLOv8n-416. These comparisons establish a multi-objective selection
problem; they do not show which model will preserve its ranking after
conversion and execution on either edge device.

The historical exact-device evidence shows that the final Sentinel TensorRT
and Scout TFLite artifacts were executable in their assigned roles. It cannot
complete the six-candidate selection question because only that deployed pair
was measured under the retained 100-iteration device protocol. The answer to
RQ1 is therefore descriptive for offline predictive trade-offs and historical
deployment, but incomplete for candidate-by-platform support.

\subsection*{RQ2: sustainable operation with concurrent services}

The ten-minute run first showed that both live cameras and deployed inference
paths could coexist with ROS~2, logging and resource sampling. The later
complete-stack run extended that observation beyond 30 minutes and added Scout
PX4 services plus Sentinel Pixhawk/MAVProxy telemetry. Both wrappers closed
normally, no logger drop or OOM occurred, minimum available memory remained
1,589.9~MiB on Sentinel and 4,873.8~MiB on Scout, and maximum valid reported
temperatures were 51.0 and 48.2~\(^{\circ}\)C.

RQ2 therefore has a bounded rather than universal answer. The evidence supports
stationary service coexistence for the measured interval, not unlimited or
airborne sustainability. Moreover, the corrected tiered-workflow probe covered only
the final 180.245~s. Its 2.618~s p95 alert-to-verdict time exceeded the
registered 2~s target, and Sentinel swap increased by 0.922~MiB. Preserving
these failures is important: resource-constrained operation was feasible under
the observed load, but the registered configuration was not compliant with
every performance criterion.

\subsection*{RQ3: temporal filtering and secondary verification}

The stationary evidence demonstrates a functional second decision stage. The
Scout confirmed all 13 requests in the stronger visible-flame interval,
confirmed one of ten requests in the smaller unstable-flame interval, and
rejected all six forced unlit requests. A fixed acknowledgement would not have
produced this condition-sensitive pattern. The event identity chain also shows
that the returned verdict was connected to the correct Sentinel request.

The evidence does not quantify improvement. The forced unlit requests bypassed
Sentinel detection, so they cannot measure false alerts prevented by the
tiered workflow. The logs did not retain every pre-filter candidate and filter-counter
transition, so they cannot reconstruct how many single-frame candidates the
three-frame rule removed or how many labelled hazards it retained. A
single-frame, filtered and verified comparison on a shared labelled sequence is
still required.

The measured cost of verification is material. Mean Scout inference exceeded
1.25~s in both accepted positive trials, and mean Sentinel
alert-to-final-decision time was approximately 1.44--1.46~s. The architecture
therefore adds an independent RGB decision at the cost of delayed resolution.
Whether that trade-off is appropriate depends on the operational cost of an
unsupported alert relative to the cost of waiting for confirmation. The
present study measures the latency but does not define a safety-critical
deadline for an industrial plant.

The repeated fault cases strengthen the coordination result without turning it
into a reliability estimate. Camera-interface loss and a post-preflight stale
stream failed closed in all five registered repetitions; absence of the Scout
process generated one timeout per event; and all ten verdicts injected after
timeout were rejected without an additional final decision. These outcomes
demonstrate specific defensive paths. They do not measure natural fault rates,
wireless recovery or rejection-memory effectiveness.

\subsection*{RQ4: two-vehicle mission behaviour and faults}

The simulation pilots demonstrate executable nominal branches and selected
safe-state transitions through real PX4 SITL and namespaced ROS~2 interfaces.
They also exposed a genuine state-assignment defect: two pre-fix fault pilots
assigned \texttt{TAKEOFF} before Scout actuation. The corrected v2 runs moved
directly from \texttt{PREFLIGHT} to \texttt{DEGRADED}. Retaining both versions
is stronger engineering evidence than reporting only the repaired runs because
it shows the failure, correction and regression outcome.

Six passing pilots do not support a reliability estimate. The rejected nominal
scenario required six mode/arm attempts after temporary health failures, which
shows that a final pass can contain operational variability. The incomplete
310-run campaign, provisional airframe proxies and scenario-defined target
coordinates prevent a broader claim about mission correctness. RQ4 is therefore
answered for the existence of the implemented branches, not for their
probability of success or their physical-flight behaviour.

\subsection*{RQ5: external transfer}

RQ5 has no evaluated performance answer. FASDD-UAV had not been acquired and locally
audited. MmodalFire passed archive and decode checks, but only 8 of 28 video
ranges passed the stricter timestamp consolidation and the event labels and
overlap gates were still open. Running the detector before those gates closed
would have produced numbers without a defensible denominator or independence
claim. Reporting the question as not evaluated preserves the value of the future
evaluation and avoids using readiness work as surrogate accuracy evidence.

\section{Predictive performance and resource trade-offs}
\label{sec:discussion:2}

The comparison illustrates why predictive performance, runtime portability and
device feasibility must remain separate selection dimensions. An offline
checkpoint can rank well and still lack a faithful conversion. The two
416-pixel YOLO ONNX artifacts passed the frozen fidelity gate, while both
640-pixel ONNX exports and the SSD TFLite export failed. Successful export and
runtime allocation were therefore insufficient indicators of equivalent
behaviour.

The YOLOv8n-640 ONNX result is a useful borderline example. Its match rate of
0.9496 was close to the 0.95 rule, and its matched confidence error was small.
Rounding the result upward or relaxing the gate would nonetheless change the
experiment after observing the outcome. Preserving the failure supports a
repeatable comparison and leaves room for a separately registered alternative
conversion.

The historical deployed pair reflects system roles rather than a model-family
optimum. The Sentinel 640-pixel TensorRT pipeline averaged 58.151~ms in the
camera-free benchmark, while the Scout TFLite pipeline averaged 999.431~ms.
The devices, models, runtimes and output interfaces all differed. The ratio
supports the observation that the deployed Sentinel path was faster than the
deployed Scout path; it does not isolate whether hardware, model size, runtime,
delegate use or implementation caused the difference.

The decoder experiment reinforces this point. Vectorization reduced the
historical Sentinel complete-pipeline mean from 204.324 to 57.049~ms while
backend inference changed only slightly. Before optimization, Python
post-processing hid much of the benefit of TensorRT. On the Scout, inference
already dominated, so the same decoder change yielded a smaller pipeline
reduction. Resource-aware design therefore requires profiling the complete
software path on each platform rather than inferring the bottleneck from model
architecture alone.

No energy-efficiency ranking is supported. Laptop resource samples combined
host-wide and device-wide quantities over unequal durations, while the physical
complete-stack record did not include calibrated power measurements. The
30-minute temperature and resource series supports bounded thermal and memory
observations, but not a causal throttling or energy-efficiency conclusion. The
strengthened device screen remains necessary before all candidate models can be
compared on common platform-specific objectives.

\section{Contribution and cost of the tiered workflow}
\label{sec:discussion:3}

The primary architectural contribution is the explicit division of
responsibilities. The Sentinel performs initial hazard detection, persistence
filtering and request publication. The Scout acquires a current frame, makes a
second model-based judgement and returns a correlated verdict. The Sentinel
then records the final decision. Correlation fields, typed outcomes and
asynchronous logging turn two unrelated detector calls into an auditable
decision workflow.

This division provides an additional rejection opportunity, but the Scout is
not independent physical ground truth. Both nodes use RGB imagery and learned
fire detectors, so their errors may be correlated through appearance, lighting,
domain shift or shared training sources. A thermal, particulate or gas sensor
would provide a different physical measurement; it was not part of the tested
verification path.

The strong- and small-flame observations show why a second decision must not be
assumed beneficial in every condition. The stronger exposure produced 13
confirmations from 13 marked requests, while the smaller unstable exposure
produced one confirmation from ten. Possible explanations include apparent
flame area, motion, viewpoint, focus, exposure, threshold calibration and
training-domain mismatch. The trial does not isolate their causal effects. It
does establish that the Scout can reject requests that the Sentinel has already
promoted, including requests associated with a real but difficult visual
target.

The communication time was small relative to Scout inference in the bounded
campaigns. Typical camera-free round-trip medians were approximately 5--6~ms,
and the loaded 1,000-event median was 4.393~ms, whereas
Scout inference was approximately 1.25--1.32~s in the accepted live conditions.
The earlier 400.453~ms communication tail remains operationally relevant:
timeouts based only on a median would ignore observed variability. Under
loaded, moving or interfered conditions, the communication share could change,
so the present ratio is not generalized beyond the retained stationary setup.

The tiered workflow also creates new failure surfaces: request loss, unavailable Scout,
stale or duplicate verdicts, incompatible clock interpretations, timeout
selection, and disagreement handling. The architecture addresses several of
these with event identity and explicit terminal states, but it does not remove
them. Its contribution is a testable and traceable coordination boundary, not
automatic reliability from adding a second vehicle.

\section{Onboard execution and warehouse-scenario implications}
\label{sec:discussion:onboard}

The implemented perception and decision path is onboard even though the
physical experiment is not aerial. The Sentinel performs capture, TensorRT
inference, filtering, coordination and finalization on its Jetson Nano, while
the Scout performs camera acquisition and TFLite verification on VOXL~2. The
Sentinel-hosted SwarmNet access point carries compact ROS~2 records between
them. The laboratory laptop supports setup, observation and evidence transfer;
it is not the required inference or final-decision computer. This distinction
is central to the title: the work evaluates distributed edge computation on
UAV hardware, but physical flight remains outside the accepted evidence.

Onboard processing has four principal advantages for the intended scenario.
First, it avoids continuous raw-video transmission and reduces the radio load
to alerts, verdicts and final decisions. Second, perception does not require an
external Internet or cloud service. Third, imagery can remain on the vehicle,
which may simplify privacy and site-data handling. Fourth, different runtimes
can be optimized for the two heterogeneous computers. The observed SwarmNet
round-trip times show that compact decision messages were inexpensive relative
to Scout inference in the stationary tests; they do not establish the same
relationship under flight or industrial radio interference.

The same design creates costs. Compute, camera and radio load draw from the
vehicle energy budget; the added Sentinel computer, regulator, camera and
mounting contribute payload; and two embedded environments must be deployed,
monitored and recovered. No calibrated electrical-power or flight-endurance
measurement quantifies those costs in this study. The Sentinel is also both
the access-point host and the final-decision coordinator. Loss of that node or
its wireless service can remove the coordination path even if the Scout can
still execute local inference. The architecture is therefore distributed in
computation but not free of concentrated dependencies.

A warehouse-like deployment adds conditions not reproduced by the laboratory.
Metallic structures and machinery can produce occlusion and multipath radio
propagation; aisles constrain line of sight and safe trajectories; illumination
can change abruptly; and smoke can obscure the RGB appearance on which both
detectors depend. Indoor GPS denial would require a validated local estimator,
shared frame convention, camera calibration, obstacle sensing and safe-standoff
planner. The simulator exercised scenario-defined local coordinates, not this
physical localization chain. Local onboard processing is consequently a
reasonable architectural response to constrained connectivity, but the present
evidence does not prove coverage, collision safety or hazard-detection
reliability in a warehouse.

\section{Failures and alternative explanations}
\label{sec:discussion:4}

The evidence contains failures at data, model, conversion, software,
integration and mission levels. Treating them as part of the result clarifies
which changes improved the system and which questions remain uncertain.

The validation-to-test rank change may arise from domain composition rather
than training instability alone. The source-family split reduces direct
leakage, but reconstructed identities cannot capture every related frame or
scene. Annotation judgement, retained negatives and the scarcity of
COCO-defined small boxes can also change metrics. One seed per configuration
does not separate model choice from training variance.

Conversion failures have several plausible sources: exporter graph changes,
preprocessing differences, coordinate conventions, suppression details and
runtime kernels. The Scout live bring-up supplied a concrete instance. The
camera stream required explicit UYVY decoding, and the model output used
normalized coordinates while the shared decoder initially expected model-pixel
coordinates. Those defects produced repeated zero-confidence outcomes until
the interfaces were inspected and corrected. The early records are therefore
configuration failures, not evidence that the detector could not see flame.

The 0.05 Scout threshold was selected after observing diagnostic illuminated
and unlit scores. It made the later proof-of-function trials possible, but it
was not preregistered or calibrated on a representative Scout validation set.
This choice can explain part of the observed confirmation pattern and prevents
the live counts from becoming an unbiased accuracy estimate.

Hardware state also affected execution. The Scout battery depleted and the
device rebooted, clearing volatile deployment files. Three of five repeated
negative JPEG payloads were lost after their hashes had been logged, although
the structured event records remained complete. The incident supports the need
for persistent staging and transfer checks. It does not invalidate the retained
scores, but it limits later visual inspection of those three frames.

The simulation pre-fix state error and temporary arming rejections show two
different failure classes. The first was a controller defect corrected in code.
The second came from PX4 health readiness and required no bypass. Combining
them under a single mission-failure count would conceal whether the remedy is a
software change, a preflight wait policy or a simulator configuration update.

\section{Threats to validity}
\label{sec:discussion:5}

\subsection{Construct validity}

The detector operationalizes visible flame appearance in RGB images. It does
not measure heat, combustion products or particulate concentration. A fire
bounding box is therefore not proof of a physical fire. Physical smoke was not
evaluated in the final two-device campaign and is outside the validated hazard
scope.

Hazard verification is also limited to a second RGB model judgement. The final
alarm was a logical decision and retained log output, not an independently
certified emergency alarm. Image-space coordinates were transferred as event
metadata, but the Scout used full-frame inference and did not use those
coordinates to crop, navigate or localize a target.

\subsection{Internal validity}

The physical trials were stationary, with propellers removed and a handheld
lighter approximately 20--25~cm from the cameras. Camera height, flame
stability, focus, exposure, viewpoint and battery state were not fully
controlled. The accepted positive intervals were marked independently, but
requests within one interval were temporally dependent. Historical
unsynchronized lighter trials cannot supply ignition-to-detection latency or
frame-level false-negative counts.

Framework-specific preprocessing and decoding could not be made identical for
all six candidates. The fidelity programme measures some of that risk but does
not remove native-framework differences. Laptop runs also occurred
sequentially and may include operating-system or display workload. These
conditions restrict causal statements about architecture alone.

\subsection{External validity}

The curated internal images and one laboratory room do not represent the
range of industrial sites, weather, lighting, occlusion, smoke, vibration,
altitude, radio interference or fire scale. The lighter stimulus is especially
different from an industrial hazard in fuel, geometry, heat and background
context. The term industrial in the title identifies the application problem,
not the validated deployment environment.

The development simulator used provisional airframe proxies and
scenario-defined local-ENU targets. Its successful state transitions do not
establish real estimator behaviour, aerodynamics, collision clearance or
radio performance. Physical two-vehicle flight generalization remains outside
the evidence.

\subsection{Statistical conclusion validity}

All six training configurations used one controlled seed. The grouped
bootstrap resampled inferred source families and therefore reflects only
within-dataset resampling variability under that grouping. It does not correct
selection bias, label error, unknown dependence or between-seed variance.
Overlapping intervals are not treated as a formal equivalence or difference
test.

The live request counts are small and dependent within exposures. Fractions
such as 13/13 and 1/10 are direct condition-specific observations rather than
binomial estimates of population recall. The 300-event communication campaigns
use serial messages from consecutive sessions, and 300/300 replies do not prove
zero DDS packet loss because retransmission can occur below the application
layer.

\begin{table}[H]
  \centering
  \caption{Principal validity threats, their effect on interpretation and the
  validation needed to reduce them.}
  \label{tab:validity-followup}
  \scriptsize
  \begin{tabularx}{\textwidth}{@{}p{0.21\textwidth}p{0.34\textwidth}X@{}}
    \toprule
    Threat & Consequence for this thesis & Required follow-up \\
    \midrule
    RGB construct only & Verification does not establish physical combustion
    or smoke presence. & Add calibrated thermal, particulate or gas sensing
    under an approved safety protocol. \\
    Small, stationary lighter trials & Counts do not estimate industrial or
    flight reliability. & Preregister repeated conditions, synchronized
    labels, varied target scale, distance, viewpoint and hard negatives. \\
    One seed per model & Ranking may include training-run variance. & Repeat
    selected candidates over multiple frozen seeds. \\
    Inferred source grouping & Residual relatedness and domain bias may remain.
    & Strengthen provenance and collect independently sourced evaluation data. \\
    Incomplete device matrix & Platform-specific model support is unknown. &
    Complete fidelity-valid 20/200 cells on both devices. \\
    Single complete-stack endurance run & Repeatability, calibrated power and
    airborne thermal behaviour remain unknown. & Repeat the registered run
    across controlled ambient and power conditions, then add staged flight only
    under an approved safety protocol. \\
    Missing stage ablation & Filter and tiered-workflow effect sizes are unknown. &
    Replay the same labelled sequence through all three decision stages. \\
    Development simulation only & No mission success probability or physical
    flight claim is supported. Qualitative review of the visual delivery checks
    presentation clarity, not simulator fidelity. & Freeze airframe parameters
    and execute the formal retained-run matrix before staged physical tests. \\
    External manifests blocked & Transfer behaviour was not evaluated. & Close
    license, alignment, overlap and eligibility gates before inference. \\
    \bottomrule
  \end{tabularx}
\end{table}

\section{Remaining validation and deployment boundaries}
\label{sec:discussion:6}

The current contribution ends at a stationary two-node perception workflow and
development simulation. The Sentinel Pixhawk participated only through
receive-only telemetry in the accepted complete-stack endurance run; no
integrated arming, take-off, navigation, collision
avoidance, return or landing was demonstrated on the real vehicles. The word
swarm denotes coordinated heterogeneous roles in a two-node architecture; it
does not imply a validated autonomous formation.

Localization is likewise incomplete. The Sentinel message carries image-space
coordinates and event context, while the simulator uses scenario-defined local
truth. A physical system would need camera calibration, a transform from image
and vehicle frames to a shared world frame, estimator uncertainty and a
validated safe standoff trajectory. Without those elements, the architecture
supports target designation at the messaging level but not geographic
localization or autonomous Scout dispatch.

Operational deployment would require more than additional accuracy data. It
would need hazard-specific safety requirements, alarm deadlines, human
oversight, fail-safe behaviour, data-governance rules, radio planning,
environmental qualification and compliance with the applicable aviation and
site procedures. Those requirements were not evaluated and are not inferred
from the laboratory or simulator results.

The evidence supports the central hypothesis to a limited degree: roles can be
divided across the tested heterogeneous edge nodes, and the implemented
workflow can produce correlated confirmations and rejections. Evidence is
insufficient to conclude that the division is sustainable beyond the measured
stationary interval, improves accuracy relative to a common baseline, transfers
to eligible external domains or operates correctly as a physical two-vehicle
industrial mission. Chapter
\ref{ch:conclusion} therefore states a conditional conclusion and prioritizes
the missing validations rather than presenting the prototype as a completed
deployment.

\reviewstatus{The discussion answers RQ1, RQ2, RQ3 and RQ4 only within their
completed evidence layers. RQ2 has a bounded answer and RQ5 was not evaluated. No
result establishes physical smoke detection, geographic localization,
autonomous industrial flight or universal reliability.}
